% =====================================================================
%  CHAPITRE 8 — GRAMMAIRE A1 : LES FONDATIONS
% =====================================================================
\chapter{Les fondations : pronoms, \en{be}, \en{have}, présent simple}

\begin{quote}\itshape
Vous n'avez pas besoin de connaître toute la grammaire pour parler. Vous avez besoin de
connaître \textbf{parfaitement} les dix structures qui représentent 70\,\% des phrases
réelles. Ce chapitre et le suivant vous les donnent.
\end{quote}

Chaque unité suit le même format : \textbf{règle} \ra \textbf{tableau} \ra \textbf{exemples
bilingues} \ra \textbf{piège francophone} \ra \textbf{exercices corrigés} \ra
\textbf{mission orale}. Ne sautez jamais la mission orale.

\section{Unité 1 --- Les pronoms personnels et le verbe \en{to be}}

\subsection{Les pronoms personnels}
\begin{tableau}[Pronoms personnels sujets et compléments]
\begin{tabularx}{\linewidth}{@{}p{3.2cm}p{3.2cm}p{3.0cm}X@{}}
\entete{Sujet} & \entete{Complément} & \entete{Possessif} & \entete{Exemple} \tabularnewline
\hline
\en{I} \pron{aï} --- je & \en{me} --- me, moi & \en{my / mine} & \en{I love you. She loves me.} \tabularnewline
\en{you} --- tu, vous & \en{you} --- te, toi, vous & \en{your / yours} & \en{You are late. I call you.} \tabularnewline
\en{he} --- il & \en{him} --- le, lui & \en{his / his} & \en{He is my brother. I see him.} \tabularnewline
\en{she} --- elle & \en{her} --- la, elle, lui & \en{her / hers} & \en{She is a doctor. I know her.} \tabularnewline
\en{it} --- il, elle, ça (choses/animaux) & \en{it} --- le, ça & \en{its / ---} & \en{It is cold. I like it.} \tabularnewline
\en{we} --- nous & \en{us} --- nous & \en{our / ours} & \en{We are ready. Help us.} \tabularnewline
\en{you} --- vous & \en{you} --- vous & \en{your / yours} & \en{You are my friends.} \tabularnewline
\en{they} --- ils, elles & \en{them} --- les, eux, elles & \en{their / theirs} & \en{They are here. I told them.} \tabularnewline
\end{tabularx}
\end{tableau}

\begin{piege}[L'anglais n'a pas de vouvoiement --- et c'est une bonne nouvelle]
\en{you} veut dire \textbf{tu} et \textbf{vous}, au singulier et au pluriel. Pas d'hésitation,
pas de faute possible. En revanche, on s'adresse presque toujours aux gens par leur prénom,
même au travail. Dire \en{Mr Diop} à un collègue peut paraître distant ; \en{Moussa} est normal.
\end{piege}

\begin{piege}[\en{it} : le pronom obligatoire]
Le français dit « il pleut », « c'est froid », « il est 8 heures ». L'anglais exige toujours
un sujet explicite : \en{it rains}, \en{it is cold}, \en{it is eight o'clock}.
\textbf{Ne dites jamais} \en{*Is cold} ou \en{*Rains}. C'est la faute n°1 des francophones
débutants.
\end{piege}

\subsection{Le verbe \en{to be} (être)}
\begin{tableau}[\en{to be} au présent --- avec contractions]
\begin{tabularx}{\linewidth}{@{}p{2.6cm}p{3.4cm}p{3.4cm}X@{}}
\entete{Forme} & \entete{Complète} & \entete{Contractée} & \entete{Traduction} \tabularnewline
\hline
Affirmation & \en{I am} & \en{I'm} \pron{aïm} & je suis \tabularnewline
 & \en{you are} & \en{you're} \pron{iou-ar} & tu es / vous êtes \tabularnewline
 & \en{he / she / it is} & \en{he's / she's / it's} & il / elle / c'est \tabularnewline
 & \en{we are} & \en{we're} \pron{ouii-ar} & nous sommes \tabularnewline
 & \en{they are} & \en{they're} \pron{dhè-ar} & ils / elles sont \tabularnewline
\hline
Négation & \en{I am not} & \en{I'm not} & je ne suis pas \tabularnewline
 & \en{you are not} & \en{you aren't} / \en{you're not} & tu n'es pas \tabularnewline
 & \en{he is not} & \en{he isn't} / \en{he's not} & il n'est pas \tabularnewline
\hline
Question & \en{Am I\ldots?} & --- & suis-je\ldots ? \tabularnewline
 & \en{Are you\ldots?} & --- & es-tu\ldots ? \tabularnewline
 & \en{Is he\ldots?} & --- & est-il\ldots ? \tabularnewline
\end{tabularx}
\end{tableau}

\begin{astuce}[À l'oral, utilisez toujours les formes contractées]
\en{I am French} sonne comme un robot ou comme une insistance (\textit{MOI, je suis français !}).
La forme naturelle est \en{I'm French.} En conversation, les contractées représentent 95\,\%
des cas. Les formes complètes s'utilisent en fin de phrase (\en{Yes, I am.}) et pour insister.
\end{astuce}

\subsection{Les six emplois de \en{be} que le français rend autrement}
\begin{tableau}[\en{be} là où le français utilise \en{avoir}]
\begin{tabularx}{\linewidth}{@{}p{6.6cm}X@{}}
\entete{Anglais} & \entete{Français} \tabularnewline
\hline
\en{I \textbf{am} 25 years old.} & J'\textbf{ai} 25 ans. \tabularnewline
\en{She \textbf{is} hungry / thirsty / cold.} & Elle \textbf{a} faim / soif / froid. \tabularnewline
\en{He \textbf{is} afraid of dogs.} & Il \textbf{a} peur des chiens. \tabularnewline
\en{We \textbf{are} late.} & Nous \textbf{sommes} en retard. (identique) \tabularnewline
\en{They \textbf{are} right.} & Ils \textbf{ont} raison. \tabularnewline
\en{How old \textbf{are} you?} & Quel âge \textbf{as}-tu ? \tabularnewline
\end{tabularx}
\end{tableau}

\begin{piege}[\en{*I have 25 years}]
Erreur systématique. Retenez le bloc entier : \en{I am X years old} --- ou, en conversation
courante, simplement \en{I'm 25}. De même : \en{*I have cold} \ra \en{I'm cold} ;
\en{*I have fear} \ra \en{I'm afraid}.
\end{piege}

\begin{exo}[1.1 --- Complétez avec \en{be}]
\begin{enumerate}
  \item I \trou[1.2cm] from Senegal. \item She \trou[1.2cm] my sister.
  \item They \trou[1.2cm] not at home. \item \trou[1.2cm] you married?
  \item We \trou[1.2cm] hungry. \item He \trou[1.2cm] 30 years old.
  \item It \trou[1.2cm] very hot today. \item Where \trou[1.2cm] your parents?
  \item I \trou[1.2cm] tired but I \trou[1.2cm] happy. \item \trou[1.2cm] she a doctor? --- No, she \trou[1.2cm].
\end{enumerate}
\end{exo}
\begin{corrige}[1.1]
1. \en{am} ('m) \quad 2. \en{is} ('s) \quad 3. \en{are} \quad 4. \en{Are} \quad
5. \en{are} \quad 6. \en{is} \quad 7. \en{is} \quad 8. \en{are} \quad
9. \en{am} / \en{am} \quad 10. \en{Is} --- \en{she isn't}
\end{corrige}

\begin{exo}[1.2 --- Traduisez en anglais]
\begin{enumerate}
  \item J'ai 32 ans. \item Elle a faim. \item Nous avons raison.
  \item Il a peur de la mer. \item Ils ne sont pas prêts. \item Quel âge as-tu ?
  \item C'est mon ami. \item Je ne suis pas fatigué.
\end{enumerate}
\end{exo}
\begin{corrige}[1.2]
1. \en{I'm 32 (years old).} 2. \en{She's hungry.} 3. \en{We're right.}
4. \en{He's afraid of the sea.} 5. \en{They aren't ready.} 6. \en{How old are you?}
7. \en{He's / This is my friend.} 8. \en{I'm not tired.}
\end{corrige}

\begin{mission}[Unité 1]
Dites à voix haute, en vous regardant dans un miroir, 10 fois de suite, en variant le
pronom : \en{I'm happy --- you're happy --- he's happy --- she's happy --- it's good ---
we're happy --- they're happy}. Puis présentez-vous en 5 phrases :
\en{Hello, my name is\ldots\ I'm from\ldots\ I'm \ldots years old. I'm a \ldots\ I'm not
married / I'm married.}
\end{mission}


\section{Unité 2 --- \en{have}, \en{have got} et la possession}

\begin{regle}[\en{have} = avoir]
\en{I / you / we / they have} \quad---\quad \en{he / she / it has}\\
Contractions : \en{I've, you've, we've, they've, he's, she's, it's}.\\
Négation : \en{I don't have} / \en{he doesn't have} (américain courant) \quad ou \quad
\en{I haven't got} / \en{he hasn't got} (britannique).\\
Question : \en{Do you have\ldots?} / \en{Have you got\ldots?}
\end{regle}

\begin{tableau}[\en{have} ou \en{have got} ?]
\begin{tabularx}{\linewidth}{@{}p{4.0cm}p{5.2cm}X@{}}
\entete{} & \entete{Américain (recommandé)} & \entete{Britannique} \tabularnewline
\hline
J'ai une voiture. & \en{I have a car.} / \en{I've got a car.} & \en{I've got a car.} \tabularnewline
Tu as un stylo ? & \en{Do you have a pen?} & \en{Have you got a pen?} \tabularnewline
Il n'a pas d'argent. & \en{He doesn't have any money.} & \en{He hasn't got any money.} \tabularnewline
\end{tabularx}
\end{tableau}

\begin{piege}[\en{He doesn't \textbf{has}}]
Après \en{does / doesn't / did / didn't}, le verbe revient à sa forme \textbf{de base} :
\en{He doesn't have} (et non \en{*doesn't has}). Même chose aux questions :
\en{Does she have a car?} (et non \en{*Does she has}).
\end{piege}

\subsection{Adjectifs et pronoms possessifs}
\begin{tableau}[Les possessifs]
\begin{tabularx}{\linewidth}{@{}p{2.6cm}p{3.4cm}p{3.4cm}X@{}}
\entete{Personne} & \entete{Adjectif + nom} & \entete{Pronom (seul)} & \entete{Exemple} \tabularnewline
\hline
je & \en{my} --- mon & \en{mine} --- le mien & \en{my book / the book is mine} \tabularnewline
tu / vous & \en{your} --- ton, votre & \en{yours} & \en{your name / it's yours} \tabularnewline
il & \en{his} --- son (à lui) & \en{his} & \en{his car} \tabularnewline
elle & \en{her} --- son (à elle) & \en{hers} & \en{her bag} \tabularnewline
il/elle (chose) & \en{its} --- son (à cela) & --- & \en{its colour} \tabularnewline
nous & \en{our} --- notre & \en{ours} & \en{our house} \tabularnewline
ils/elles & \en{their} --- leur & \en{theirs} & \en{their children} \tabularnewline
\end{tabularx}
\end{tableau}

\begin{memo}[Le possessif suit le \textbf{possesseur}, pas l'objet]
En français, « son » dépend du genre de l'objet possédé (\emph{son} livre, \emph{sa} voiture).
En anglais, il dépend de \textbf{qui possède} : \en{his book} (le livre d'un homme),
\en{her book} (le livre d'une femme) --- même si c'est le même objet. Et \en{his} ou \en{her}
ne changent jamais au pluriel : \en{his books}, \en{her cars}.
\end{memo}

\subsection{Le génitif : \en{'s} (appartenance)}
\begin{regle}
Nom + \en{'s} = « de » : \en{Moussa's car} (la voiture de Moussa), \en{my brother's wife}
(la femme de mon frère), \en{the children's room}. Au pluriel en \en{-s}, on met seulement
l'apostrophe : \en{the students' books}. Pour les choses, on préfère \en{of} :
\en{the door of the car}, \en{the end of the film}.
\end{regle}

\begin{exo}[2.1 --- Traduisez]
\begin{enumerate}
  \item J'ai deux frères. \item Elle a une belle maison. \item Nous n'avons pas de temps.
  \item Est-ce qu'il a une voiture ? \item Le téléphone de Fatou. \item Le sac des enfants.
  \item C'est le mien. \item Leur maison est grande.
\end{enumerate}
\end{exo}
\begin{corrige}[2.1]
1. \en{I have two brothers.} 2. \en{She has a beautiful house.}
3. \en{We don't have (any) time.} 4. \en{Does he have a car?} / \en{Has he got a car?}
5. \en{Fatou's phone.} 6. \en{The children's bag.} 7. \en{It's mine.}
8. \en{Their house is big.}
\end{corrige}

\section{Unité 3 --- Le présent simple}

\begin{regle}[Formation]
\textbf{Toutes les personnes : base verbale}, sauf la 3\textsuperscript{e} personne du
singulier qui prend \en{-s}.\\
\en{I work, you work, he work\textbf{s}, she work\textbf{s}, it work\textbf{s}, we work,
they work.}\\
Négation : \en{do not (don't)} + base --- \en{I don't work}, \en{he doesn't work}.\\
Question : \en{do / does} + sujet + base --- \en{Do you work?}, \en{Does she work?}
\end{regle}

\begin{tableau}[Le \en{-s} de la 3\textsuperscript{e} personne : orthographe]
\begin{tabularx}{\linewidth}{@{}p{5.6cm}X@{}}
\entete{Cas} & \entete{Règle et exemples} \tabularnewline
\hline
Règle générale & + \en{s} : \en{works, plays, reads, eats} \tabularnewline
Verbes en \en{-o}, \en{-ch}, \en{-sh}, \en{-ss}, \en{-x} & + \en{es} : \en{goes, watches, washes, passes, fixes} \tabularnewline
Consonne + \en{y} & \en{y} \ra \en{ies} : \en{study} \ra \en{studies}, \en{try} \ra \en{tries}, \en{cry} \ra \en{cries} \tabularnewline
Voyelle + \en{y} & + \en{s} simplement : \en{plays, buys, says, enjoys} \tabularnewline
Irréguliers & \en{have} \ra \en{has} \tabularnewline
\end{tabularx}
\end{tableau}

\subsection{Quand employer le présent simple ?}
\begin{enumerate}
  \item \textbf{Habitudes et routines} : \en{I wake up at six every day.}
  \item \textbf{Vérités générales} : \en{Water boils at 100 degrees.} \en{The sun rises in the east.}
  \item \textbf{États permanents} : \en{She lives in Thiès.} \en{He works in a bank.}
  \item \textbf{Horaires officiels} (futur) : \en{The train leaves at 8 pm.}
  \item \textbf{Avec les verbes d'état} (voir, entendre, aimer, savoir\ldots) :
  \en{I know, I understand, I want, I need, I like, I love, I prefer, I believe, I remember}.
\end{enumerate}

\begin{piege}[Les verbes d'état ne se mettent pas au \en{-ing}]
\en{*I am knowing the answer} \ra \en{I know the answer}.
\en{*I am wanting a coffee} \ra \en{I want a coffee}.
Liste à connaître : \en{know, want, need, like, love, hate, prefer, believe, understand,
remember, forget, mean, belong, own, seem, sound, smell, taste, feel (sens), hope, wish}.
\end{piege}

\subsection{Les adverbes de fréquence et leur place}
\begin{tableau}[Échelle de fréquence + placement]
\begin{tabularx}{\linewidth}{@{}p{3.4cm}p{2.0cm}X@{}}
\entete{Adverbe} & \entete{\%} & \entete{Exemple (placement)} \tabularnewline
\hline
\en{always} & 100 & \en{I \textbf{always} drink tea.} \tabularnewline
\en{usually} & 90 & \en{She \textbf{usually} takes the bus.} \tabularnewline
\en{often} & 70 & \en{We \textbf{often} watch football.} \tabularnewline
\en{sometimes} & 50 & \en{They \textbf{sometimes} cook fish.} (aussi en début ou fin de phrase) \tabularnewline
\en{rarely / seldom} & 10 & \en{He \textbf{rarely} goes out.} \tabularnewline
\en{never} & 0 & \en{I \textbf{never} eat pork.} \tabularnewline
\en{hardly ever} & 2 & \en{We \textbf{hardly ever} travel.} \tabularnewline
\end{tabularx}
\end{tableau}

\begin{regle}[Placement]
\textbf{Avant} le verbe principal : \en{I \textbf{never} smoke.}\\
\textbf{Après} le verbe \en{be} : \en{He is \textbf{always} late.}\\
\textbf{Entre} l'auxiliaire et le verbe : \en{I have \textbf{never} been to London.}
\en{She doesn't \textbf{often} call.}
\end{regle}

\begin{piege}[Jamais de double négation]
Le français dit « je ne sais \emph{pas} » ; le wolof et d'autres langues autorisent la double
négation. L'anglais \textbf{non} : \en{never} et \en{not} ne cohabitent jamais.\\
\en{*I don't never go there.} \ra \en{I never go there.} ou \en{I don't ever go there.}
\end{piege}

\begin{exo}[3.1 --- Mettez le verbe à la bonne forme]
\begin{enumerate}
  \item She (work) \trou[2cm] in a hospital.
  \item They (not / like) \trou[2cm] coffee.
  \item (you / speak) \trou[2.5cm] English?
  \item My brother (study) \trou[2cm] every evening.
  \item He (go) \trou[2cm] to the mosque on Friday.
  \item We (watch) \trou[2cm] the news at 8 pm.
  \item (she / have) \trou[2.5cm] a car?
  \item I (not / understand) \trou[2.5cm] this word.
  \item The bus (leave) \trou[2cm] at 7 am.
  \item My parents (live) \trou[2cm] in Dakar.
\end{enumerate}
\end{exo}
\begin{corrige}[3.1]
1. \en{works} 2. \en{don't like} 3. \en{Do you speak} 4. \en{studies} 5. \en{goes}
6. \en{watch} 7. \en{Does she have} 8. \en{don't understand} 9. \en{leaves} 10. \en{live}
\end{corrige}

\begin{exo}[3.2 --- Placez l'adverbe de fréquence]
\begin{enumerate}
  \item I drink tea. (always) \item She is late. (never)
  \item We go to the beach. (sometimes) \item He watches TV. (hardly ever)
  \item They arrive on time. (usually) \item I have eaten sushi. (never)
\end{enumerate}
\end{exo}
\begin{corrige}[3.2]
1. \en{I always drink tea.} 2. \en{She is never late.} 3. \en{We sometimes go to the beach.}
4. \en{He hardly ever watches TV.} 5. \en{They usually arrive on time.}
6. \en{I have never eaten sushi.}
\end{corrige}

\begin{mission}[Unité 3 --- votre routine, 3 minutes]
Décrivez votre journée habituelle en 12 phrases au présent simple, à voix haute, avec des
adverbes de fréquence : \en{I always wake up at\ldots\ I usually drink\ldots\ I never\ldots\
Sometimes I\ldots} Enregistrez-vous. Réécoutez le lendemain et corrigez vos \en{-s} oubliés
--- c'est l'erreur la plus fréquente à l'oral.
\end{mission}


\section{Unité 4 --- Les articles : \en{a}, \en{an}, \en{the}, ou rien}

C'est \textbf{le} point où les francophones font des fautes à vie. La raison : le français
met un article presque partout, l'anglais souvent pas.

\begin{regle}[Les trois cas]
\begin{enumerate}
  \item \en{a / an} = \textbf{un, une} (singulier, non précis, première mention).
  \en{a} devant consonne \emph{sonore}, \en{an} devant voyelle \emph{sonore} :
  \en{a book, an apple, a university} (\pron{iou} = consonne), \en{an hour} (\en{h} muet).
  \item \en{the} = \textbf{le, la, les} (chose précise, déjà mentionnée, unique).
  \item \textbf{Aucun article} (article « zéro ») : généralités, pluriels et indénombrables
  pris en général, noms de pays (sauf exceptions), repas, langues, matières, sports.
\end{enumerate}
\end{regle}

\begin{tableau}[Article zéro : les 10 cas à connaître]
\begin{tabularx}{\linewidth}{@{}p{5.4cm}p{5.0cm}X@{}}
\entete{Règle} & \entete{Anglais} & \entete{Français (avec article !)} \tabularnewline
\hline
Généralité avec un pluriel & \en{Dogs are friendly.} & \textbf{Les} chiens sont gentils. \tabularnewline
Généralité avec un indénombrable & \en{Water is essential.} & \textbf{L'}eau est essentielle. \tabularnewline
Repas & \en{I had breakfast at 7.} & J'ai pris \textbf{le} petit-déjeuner. \tabularnewline
Langues & \en{She speaks English.} & Elle parle \textbf{l'}anglais. \tabularnewline
Matières & \en{This table is wood.} / \en{made of wood} & Cette table est \textbf{en} bois. \tabularnewline
Sports et jeux & \en{I play football.} & Je joue \textbf{au} football. \tabularnewline
Noms de pays (singulier sans \en{States/Kingdom/Republic}) & \en{I live in Senegal.} & Je vis \textbf{au} Sénégal. \tabularnewline
Villes & \en{She lives in Thiès.} & Elle vit \textbf{à} Thiès. \tabularnewline
Années, mois, jours & \en{in 2026, on Monday, in May} & \textbf{en} 2026, \textbf{le} lundi \tabularnewline
Institutions au sens général & \en{go to school, go to bed, go to church, go to work, go home} & aller à \textbf{l'}école, se coucher \tabularnewline
\end{tabularx}
\end{tableau}

\begin{tableau}[Les exceptions qui contredisent la règle des pays]
\begin{tabularx}{\linewidth}{@{}p{6.4cm}X@{}}
\entete{Avec \en{the}} & \entete{Exemples} \tabularnewline
\hline
Pays au pluriel ou contenant \en{States, Kingdom, Republic, Union} &
\en{the United States, the United Kingdom, the Netherlands, the Philippines,
the Democratic Republic of Congo, the European Union} \tabularnewline
Fleuves, mers, océans, déserts, chaînes de montagnes &
\en{the Senegal River, the Nile, the Atlantic, the Sahara, the Alps} \tabularnewline
Journaux, hôtels, musées, cinémas & \en{the Times, the Hilton, the Louvre} \tabularnewline
Adjectif + nationalité au pluriel & \en{the French, the British, the Senegalese} \tabularnewline
\end{tabularx}
\end{tableau}

\begin{piege}[\en{I like the music} ou \en{I like music} ?]
\en{I like music.} = j'aime la musique en général (le bon choix 9 fois sur 10).\\
\en{I like the music} = j'aime \emph{cette} musique-là (celle qu'on entend maintenant).\\
Même logique : \en{Life is beautiful} (la vie en général) / \en{The life of Mandela}
(la vie d'une personne précise). Les francophones ajoutent \en{the} par réflexe : c'est
la faute la plus tenace de toute la langue.
\end{piege}

\begin{exo}[4.1 --- \en{a}, \en{an}, \en{the} ou rien ?]
\begin{enumerate}
  \item I bought \trou[1cm] new phone yesterday. \trou[1cm] phone is blue.
  \item She is \trou[1cm] engineer. She works in \trou[1cm] Dakar.
  \item \trou[1cm] life is short.
  \item We had \trou[1cm] lunch at \trou[1cm] one o'clock.
  \item He speaks \trou[1cm] French and \trou[1cm] English.
  \item \trou[1cm] sun rises in \trou[1cm] east.
  \item I go to \trou[1cm] work by \trou[1cm] bus.
  \item \trou[1cm] United States is a big country.
  \item This is \trou[1cm] hour of my life I will never forget.
  \item She studies at \trou[1cm] university in \trou[1cm] Senegal.
\end{enumerate}
\end{exo}
\begin{corrige}[4.1]
1. \en{a} --- \en{The} \quad 2. \en{an} --- rien \quad 3. rien \quad 4. rien --- rien \quad
5. rien --- rien \quad 6. \en{The} --- \en{the} \quad 7. rien --- rien \quad
8. \en{The} \quad 9. \en{an} \quad 10. rien (ou \en{a} si on précise laquelle) --- rien
\end{corrige}

\section{Unité 5 --- Le pluriel, les dénombrables et \en{some / any}}

\begin{tableau}[Formation du pluriel]
\begin{tabularx}{\linewidth}{@{}p{5.0cm}X@{}}
\entete{Cas} & \entete{Exemples} \tabularnewline
\hline
Règle générale : + \en{s} & \en{book} \ra \en{books}, \en{car} \ra \en{cars} \tabularnewline
\en{-s, -sh, -ch, -x, -o} : + \en{es} & \en{bus} \ra \en{buses}, \en{box} \ra \en{boxes}, \en{watch} \ra \en{watches}, \en{tomato} \ra \en{tomatoes}, \en{potato} \ra \en{potatoes} (mais \en{photo} \ra \en{photos}, \en{video} \ra \en{videos}) \tabularnewline
Consonne + \en{y} : \en{y} \ra \en{ies} & \en{city} \ra \en{cities}, \en{baby} \ra \en{babies}, \en{country} \ra \en{countries} \tabularnewline
\en{-f / -fe} : \ra \en{-ves} & \en{leaf} \ra \en{leaves}, \en{wife} \ra \en{wives}, \en{knife} \ra \en{knives}, \en{life} \ra \en{lives}, \en{shelf} \ra \en{shelves} (mais \en{roof} \ra \en{roofs}) \tabularnewline
\textbf{Irréguliers} & \en{man} \ra \en{men} ; \en{woman} \ra \en{women} ; \en{child} \ra \en{children} ; \en{person} \ra \en{people} ; \en{tooth} \ra \en{teeth} ; \en{foot} \ra \en{feet} ; \en{mouse} \ra \en{mice} ; \en{goose} \ra \en{geese} \tabularnewline
\textbf{Invariables} & \en{sheep, fish, deer, series, species, aircraft} \tabularnewline
\textbf{Toujours pluriel} & \en{trousers, jeans, glasses, scissors, clothes, police, people} \tabularnewline
\end{tabularx}
\end{tableau}

\subsection{Dénombrables et indénombrables}
\begin{tableau}[Les indénombrables les plus utiles]
\begin{tabularx}{\linewidth}{@{}p{4.4cm}p{4.4cm}X@{}}
\entete{Alimentation} & \entete{Abstraits} & \entete{Matières / autres} \tabularnewline
\hline
\en{bread, rice, water, milk, sugar, salt, meat, fish, fruit, butter, cheese, tea, coffee, oil} &
\en{information, advice, news, money, work, time, help, knowledge, progress, research, furniture, luggage, equipment, traffic, homework} &
\en{paper, wood, gold, iron, hair, grass, electricity, weather, rain, snow, music, English, sleep} \tabularnewline
\end{tabularx}
\end{tableau}

\begin{piege}[Ces mots sont indénombrables en anglais et dénombrables en français]
\en{*an information} \ra \en{a piece of information}\\
\en{*an advice} \ra \en{a piece of advice}\\
\en{*a news} \ra \en{a piece of news} (\en{news} est singulier : \en{The news \textbf{is} bad.})\\
\en{*a furniture} \ra \en{a piece of furniture}\\
\en{*a luggage} \ra \en{a piece of luggage / a bag}\\
\en{*a hair} (= la chevelure) \ra \en{a hair} existe mais veut dire \emph{un poil} !\\
\en{*two breads} \ra \en{two loaves of bread} --- \en{*a money} \ra \en{some money}
\end{piege}

\subsection{\en{some}, \en{any}, \en{no}, \en{much}, \en{many}, \en{a lot of}}
\begin{tableau}[Quantifieurs de base]
\begin{tabularx}{\linewidth}{@{}p{2.8cm}p{4.0cm}X@{}}
\entete{Mot} & \entete{Emploi} & \entete{Exemples} \tabularnewline
\hline
\en{some} & Phrase affirmative ; offres et demandes polies & \en{I need some water.} / \en{Would you like some tea?} / \en{Can I have some bread?} \tabularnewline
\en{any} & Négation et question & \en{I don't have any money.} / \en{Do you have any brothers?} \tabularnewline
\en{no} & Négation sans \en{not} & \en{I have no time.} (= \en{I don't have any time.}) \tabularnewline
\en{much} & Indénombrables (négation, question) & \en{How much sugar?} / \en{I don't have much time.} \tabularnewline
\en{many} & Dénombrables & \en{How many children?} / \en{There aren't many cars.} \tabularnewline
\en{a lot of / lots of} & Les deux, à l'affirmative (le plus naturel) & \en{I have a lot of friends.} / \en{She drinks a lot of water.} \tabularnewline
\en{a few / few} & Dénombrables : quelques / très peu & \en{a few friends} (positif) / \en{few friends} (presque aucun) \tabularnewline
\en{a little / little} & Indénombrables : un peu / très peu & \en{a little sugar} / \en{little hope} \tabularnewline
\end{tabularx}
\end{tableau}

\begin{exo}[5.1 --- Pluriels]
Écrivez le pluriel : \en{child, woman, foot, tooth, person, mouse, bus, city, wife, leaf,
sheep, potato, box, baby, man, fish, glass, country, knife, life}.
\end{exo}
\begin{corrige}[5.1]
\en{children, women, feet, teeth, people, mice, buses, cities, wives, leaves, sheep,
potatoes, boxes, babies, men, fish, glasses, countries, knives, lives}.
\end{corrige}

\begin{exo}[5.2 --- \en{some}, \en{any}, \en{much}, \en{many}, \en{a lot of}]
\begin{enumerate}
  \item I don't have \trou[1.5cm] time.
  \item How \trou[1.5cm] brothers do you have?
  \item There is \trou[1.5cm] milk in the fridge.
  \item Do you want \trou[1.5cm] tea?
  \item She has \trou[1.5cm] of friends.
  \item How \trou[1.5cm] money do you need?
  \item We don't have \trou[1.5cm] bread left.
  \item There are \trou[1.5cm] people in the room.
\end{enumerate}
\end{exo}
\begin{corrige}[5.2]
1. \en{any / much} 2. \en{many} 3. \en{some} 4. \en{some} (offre) 5. \en{a lot / lots}
6. \en{much} 7. \en{any / much} 8. \en{a lot of / many}
\end{corrige}

\begin{exo}[5.3 --- Corrigez l'erreur]
\begin{enumerate}
  \item \en{*Can you give me an advice?}
  \item \en{*I have many informations for you.}
  \item \en{*The news are good today.}
  \item \en{*I bought two breads and a milk.}
  \item \en{*She has a beautiful hairs.}
  \item \en{*My luggages are heavy.}
\end{enumerate}
\end{exo}
\begin{corrige}[5.3]
1. \en{Can you give me a piece of advice?} (ou \en{some advice})
2. \en{I have a lot of information for you.}
3. \en{The news is good today.}
4. \en{I bought two loaves of bread and some milk.}
5. \en{She has beautiful hair.}
6. \en{My luggage is heavy.} / \en{My bags are heavy.}
\end{corrige}


\section{Unité 6 --- Poser des questions : la mécanique complète}

\begin{regle}[L'ordre des mots : AUX + SUJET + VERBE]
L'anglais n'a pas d'inversion sujet-verbe comme le français. On place un \textbf{auxiliaire}
devant le sujet :\\
\en{You work here.} \ra \en{\textbf{Do} you work here?}\\
\en{She is tired.} \ra \en{\textbf{Is} she tired?}\\
\en{They can swim.} \ra \en{\textbf{Can} they swim?}\\
\en{He went out.} \ra \en{\textbf{Did} he go out?} (le verbe redevient base !)
\end{regle}

\begin{tableau}[Les mots interrogatifs (\en{Wh-} words)]
\begin{tabularx}{\linewidth}{@{}p{2.8cm}p{3.4cm}X@{}}
\entete{Mot} & \entete{Sens} & \entete{Exemple} \tabularnewline
\hline
\en{What} & que, quoi, quel & \en{What do you want?} --- \textit{Que veux-tu ?} \tabularnewline
\en{What\ldots like} & comment est (description) & \en{What's he like?} --- \textit{Comment est-il ? (caractère)} \tabularnewline
\en{Where} & où & \en{Where do you live?} \tabularnewline
\en{When} & quand & \en{When does the film start?} \tabularnewline
\en{Why} & pourquoi & \en{Why are you late?} --- réponse par \en{Because\ldots} \tabularnewline
\en{Who} & qui & \en{Who is that man?} / \en{Who called?} \tabularnewline
\en{Whom} & qui (complément, formel) & \en{To whom did you speak?} \tabularnewline
\en{Whose} & à qui & \en{Whose bag is this?} \tabularnewline
\en{Which} & lequel, quel (choix limité) & \en{Which colour do you prefer?} \tabularnewline
\en{How} & comment & \en{How does it work?} \tabularnewline
\en{How much} & combien (indénombrable) & \en{How much is it?} \tabularnewline
\en{How many} & combien (dénombrable) & \en{How many people came?} \tabularnewline
\en{How long} & combien de temps & \en{How long does it take?} \tabularnewline
\en{How often} & à quelle fréquence & \en{How often do you study?} \tabularnewline
\en{How far} & à quelle distance & \en{How far is the airport?} \tabularnewline
\en{How old} & quel âge & \en{How old are you?} \tabularnewline
\en{What time} & à quelle heure & \en{What time do you get up?} \tabularnewline
\end{tabularx}
\end{tableau}

\begin{piege}[Les questions \emph{sujet} : pas d'auxiliaire !]
Quand le mot interrogatif \textbf{est} le sujet, l'ordre reste normal et le verbe prend
\en{-s} :\\
\en{Who wants coffee?} (et non \en{*Who does want coffee?})\\
\en{What happened?} (et non \en{*What did happen?})\\
\en{Who called you?} \quad \en{Which bus goes to the centre?}
\end{piege}

\subsection{Les questions indirectes (la politesse)}
\begin{regle}
Devant \en{Can you tell me\ldots}, \en{Do you know\ldots}, \en{I wonder\ldots},
\en{Could you tell me\ldots}, l'ordre redevient \textbf{affirmatif} : sujet puis verbe,
et \en{do/does/did} disparaît.
\end{regle}
\ex{Where is the station? \ra\ Could you tell me where the station is?}{Où est la gare ? \ra\ Pourriez-vous me dire où est la gare ?}
\ex{What time does it start? \ra\ Do you know what time it starts?}{À quelle heure ça commence ? \ra\ Savez-vous à quelle heure ça commence ?}
\ex{Why did he leave? \ra\ I wonder why he left.}{Pourquoi est-il parti ? \ra\ Je me demande pourquoi il est parti.}

\begin{astuce}[Les formules qui font « anglais naturel »]
\en{Do you know if\ldots?} (Savez-vous si\ldots ?) \quad \en{I was wondering whether\ldots}
(Je me demandais si\ldots) \quad \en{Sorry to bother you, but\ldots} (Pardon de vous
déranger, mais\ldots) \quad \en{Would you mind telling me\ldots?} (Cela vous ennuierait-il
de me dire\ldots ?). Ces quatre formules transforment instantanément un anglais scolaire en
anglais poli d'adulte.
\end{astuce}

\subsection{Les 30 questions à connaître par c\oe ur}
\begin{tableau}[30 questions essentielles]
\footnotesize
\begin{longtable}{@{}p{6.6cm}p{7.0cm}@{}}
\entete{Anglais} & \entete{Français} \tabularnewline
\endfirsthead
\entete{Anglais} & \entete{Français} \tabularnewline
\endhead
\phr{What's your name?}{Comment vous appelez-vous ?}
\phr{Where are you from?}{D'où venez-vous ?}
\phr{Where do you live?}{Où habitez-vous ?}
\phr{How old are you?}{Quel âge avez-vous ?}
\phr{What do you do? (for a living)}{Que faites-vous dans la vie ?}
\phr{Are you married?}{Êtes-vous marié(e) ?}
\phr{Do you have children?}{Avez-vous des enfants ?}
\phr{Do you speak English?}{Parlez-vous anglais ?}
\phr{Can you repeat that, please?}{Pouvez-vous répéter, s'il vous plaît ?}
\phr{Can you speak more slowly, please?}{Pouvez-vous parler plus lentement ?}
\phr{What does ``\ldots'' mean?}{Que veut dire « \ldots » ?}
\phr{How do you say ``\ldots'' in English?}{Comment dit-on « \ldots » en anglais ?}
\phr{How do you spell that?}{Comment cela s'épelle-t-il ?}
\phr{What time is it?}{Quelle heure est-il ?}
\phr{How much is it?}{Combien ça coûte ?}
\phr{Where is the bathroom?}{Où sont les toilettes ?}
\phr{How do I get to\ldots?}{Comment aller à\ldots ?}
\phr{Is it far from here?}{C'est loin d'ici ?}
\phr{What are you doing?}{Que fais-tu ?}
\phr{What did you do yesterday?}{Qu'as-tu fait hier ?}
\phr{Where did you go?}{Où es-tu allé ?}
\phr{Why not?}{Pourquoi pas ?}
\phr{What's wrong?}{Qu'est-ce qui ne va pas ?}
\phr{What happened?}{Que s'est-il passé ?}
\phr{How was your day?}{Comment s'est passée ta journée ?}
\phr{Are you hungry?}{As-tu faim ?}
\phr{Do you mind if I\ldots?}{Cela t'ennuie si je\ldots ?}
\phr{What do you think?}{Qu'en pensez-vous ?}
\phr{Could you help me, please?}{Pourriez-vous m'aider ?}
\phr{When will you be ready?}{Quand serez-vous prêt ?}
\end{longtable}
\end{tableau}

\begin{exo}[6.1 --- Transformez en question]
\begin{enumerate}
  \item \en{You work in a bank.} \item \en{She lives in Dakar.}
  \item \en{They arrived yesterday.} \item \en{He can drive.}
  \item \en{You are thirsty.} \item \en{She wants a coffee.}
  \item \en{Someone called.} (avec \en{who}) \item \en{You need help.} (avec \en{what})
\end{enumerate}
\end{exo}
\begin{corrige}[6.1]
1. \en{Do you work in a bank?} 2. \en{Does she live in Dakar?}
3. \en{Did they arrive yesterday?} 4. \en{Can he drive?} 5. \en{Are you thirsty?}
6. \en{Does she want a coffee?} 7. \en{Who called?} 8. \en{What do you need?}
\end{corrige}

\begin{exo}[6.2 --- Questions indirectes]
Rendez ces questions polies :
\begin{enumerate}
  \item \en{Where is the bus stop?} (avec \en{Can you tell me\ldots})
  \item \en{What time does the shop close?} (avec \en{Do you know\ldots})
  \item \en{Why is he angry?} (avec \en{I wonder\ldots})
  \item \en{How much does it cost?} (avec \en{Could you tell me\ldots})
\end{enumerate}
\end{exo}
\begin{corrige}[6.2]
1. \en{Can you tell me where the bus stop is?} 2. \en{Do you know what time the shop closes?}
3. \en{I wonder why he is angry.} 4. \en{Could you tell me how much it costs?}
\end{corrige}

\section{Unité 7 --- La négation}

\begin{tableau}[Toutes les façons de nier]
\begin{tabularx}{\linewidth}{@{}p{3.6cm}p{4.6cm}X@{}}
\entete{Type} & \entete{Construction} & \entete{Exemples} \tabularnewline
\hline
Avec \en{be} & sujet + \en{be} + \en{not} & \en{I'm not tired. He isn't here.} \tabularnewline
Avec un auxiliaire/modal & auxiliaire + \en{not} & \en{I can't swim. She won't come. We haven't eaten.} \tabularnewline
Avec un verbe simple & \en{don't / doesn't / didn't} + base & \en{I don't know. He doesn't work. They didn't come.} \tabularnewline
Avec un adverbe négatif & \en{never, hardly, rarely, nothing, nobody} & \en{I never lie. Nobody called. Nothing happened.} \tabularnewline
Réponses courtes & \en{No, I don't. / No, she can't.} & \en{Do you like it?} --- \en{No, I don't.} \tabularnewline
\end{tabularx}
\end{tableau}

\begin{memo}[Les mots négatifs suffisent seuls]
\en{nothing} = \en{not anything} ; \en{nobody / no one} = \en{not anybody} ;
\en{nowhere} = \en{not anywhere} ; \en{never} = \en{not ever}.\\
Un seul par phrase : \en{I said nothing.} \textbf{ou} \en{I didn't say anything.}
Jamais les deux (\en{*I didn't say nothing}).
\end{memo}

\begin{piege}[Les réponses \en{yes} / \en{no} après une question négative]
\en{Don't you like coffee?} --- Si vous aimez le café, répondez \en{Yes, I do!} (= « si ! »).
Le français dit « si », l'anglais dit \en{yes}. Ne traduisez pas « non » littéralement :
répondez toujours selon la \textbf{réalité}, pas selon la forme de la question.
\end{piege}

\section{Unité 8 --- \en{There is / There are} et les prépositions de lieu}

\begin{regle}
\en{There is} + singulier / indénombrable \quad---\quad \en{There are} + pluriel\\
Négation : \en{There isn't / There aren't} \quad---\quad Question : \en{Is there\ldots? / Are there\ldots?}\\
Contracté : \en{There's a problem.} (\en{there're} s'écrit rarement)
\end{regle}

\begin{tableau}[Les prépositions de lieu essentielles]
\begin{tabularx}{\linewidth}{@{}p{2.6cm}p{3.6cm}X@{}}
\entete{Préposition} & \entete{Sens} & \entete{Exemple} \tabularnewline
\hline
\en{in} & dans (espace clos, ville, pays) & \en{in the box, in Dakar, in Senegal, in the car} \tabularnewline
\en{on} & sur (surface, étage, transport collectif) & \en{on the table, on the wall, on the second floor, on the bus} \tabularnewline
\en{at} & à (point précis, adresse, lieu d'activité) & \en{at the door, at home, at work, at school, at 25 Independence Avenue} \tabularnewline
\en{under} & sous & \en{under the bed} \tabularnewline
\en{next to / beside} & à côté de & \en{next to the bank} \tabularnewline
\en{between} & entre & \en{between the shop and the mosque} \tabularnewline
\en{in front of} & devant & \en{in front of the house} \tabularnewline
\en{behind} & derrière & \en{behind the car} \tabularnewline
\en{opposite} & en face de & \en{opposite the market} \tabularnewline
\en{near / close to} & près de & \en{near the beach} \tabularnewline
\en{above / below} & au-dessus / au-dessous (sans contact) & \en{above the door} \tabularnewline
\en{into / out of} & dans (mouvement) / hors de & \en{He went into the room.} \tabularnewline
\en{through} & à travers & \en{through the window, through the city} \tabularnewline
\en{across} & de l'autre côté, en traversant & \en{across the street} \tabularnewline
\en{along} & le long de & \en{along the river} \tabularnewline
\en{around} & autour de & \en{around the table} \tabularnewline
\end{tabularx}
\end{tableau}

\begin{astuce}[\en{in the car} mais \en{on the bus}]
On dit \en{in} pour les petits véhicules où l'on s'assoit (\en{in a car, in a taxi,
in a truck, in a helicopter}) et \en{on} pour ceux où l'on \emph{marche}
(\en{on a bus, on a train, on a plane, on a ship}) --- et aussi \en{on a bike, on a motorbike,
on a horse}.
\end{astuce}

\subsection{Les démonstratifs}
\begin{tableau}[\en{this / that / these / those}]
\begin{tabularx}{\linewidth}{@{}p{3.0cm}p{3.4cm}X@{}}
\entete{} & \entete{Près de moi} & \entete{Loin de moi} \tabularnewline
\hline
Singulier & \en{this} \pron{dhis} --- ceci, ce & \en{that} \pron{dhat} --- cela, ce \tabularnewline
Pluriel & \en{these} \pron{dhiz} --- ces & \en{those} \pron{dhouz} --- ces (là-bas) \tabularnewline
\end{tabularx}
\end{tableau}
\ex{This is my friend Ali. Those are my keys.}{Voici mon ami Ali. Ce sont mes clés (là-bas).}
\ex{I don't like that idea. What's this?}{Je n'aime pas cette idée. Qu'est-ce que c'est ?}

\begin{exo}[8.1 --- Traduisez]
\begin{enumerate}
  \item Il y a un problème. \item Il n'y a pas de chaises dans la salle.
  \item Y a-t-il un hôpital près d'ici ? \item Le chat est sous la table.
  \item La banque est entre la pharmacie et le marché. \item Ces livres sont à moi.
  \item Qu'est-ce que c'est ? \item Je suis au travail.
\end{enumerate}
\end{exo}
\begin{corrige}[8.1]
1. \en{There's a problem.} 2. \en{There aren't any chairs in the room.}
3. \en{Is there a hospital near here?} 4. \en{The cat is under the table.}
5. \en{The bank is between the pharmacy and the market.} 6. \en{These books are mine.}
7. \en{What is this?} / \en{What's this?} 8. \en{I'm at work.}
\end{corrige}

\section{Unité 9 --- Les adjectifs}

\begin{regle}[Trois règles d'or]
\begin{enumerate}
  \item L'adjectif se place \textbf{avant} le nom : \en{a big house} (une grande maison).
  \item L'adjectif est \textbf{invariable} : \en{a blue car, two blue cars} (jamais \en{*blues}).
  \item Avec \en{be}, l'adjectif reste après le verbe : \en{The house is big.}
\end{enumerate}
\end{regle}

\begin{tableau}[L'ordre des adjectifs : OSASCOMP]
\begin{tabularx}{\linewidth}{@{}p{2.2cm}p{3.4cm}X@{}}
\entete{Rang} & \entete{Type} & \entete{Exemple dans l'ordre} \tabularnewline
\hline
\textbf{O} & Opinion & \en{beautiful, nice, ugly, expensive, old-fashioned} \tabularnewline
\textbf{S} & Size / shape & \en{big, small, tall, round, long} \tabularnewline
\textbf{A} & Age & \en{new, old, young, ancient} \tabularnewline
\textbf{C} & Colour & \en{red, black, dark blue} \tabularnewline
\textbf{O} & Origin & \en{Senegalese, French, Chinese} \tabularnewline
\textbf{M} & Material & \en{wooden, metal, plastic, cotton} \tabularnewline
\textbf{P} & Purpose & \en{sleeping bag, running shoes, dining room} \tabularnewline
\end{tabularx}
\end{tableau}
\ex{a beautiful big old brown Senegalese wooden house}{une belle grande vieille maison sénégalaise en bois marron}
\textit{(En pratique, on met rarement plus de trois adjectifs.)}

\subsection{Adjectifs en \en{-ed} et \en{-ing}}
\begin{regle}
\en{-ed} = ce que \textbf{je ressens} (\emph{é} comme « je suis fatigué »).\\
\en{-ing} = ce qui \textbf{cause} le sentiment (\emph{ant} comme « c'est fatigant »).
\end{regle}
\ex{I'm bored. --- The film is boring.}{Je m'ennuie. --- Le film est ennuyeux.}
\ex{She's interested in history. --- History is interesting.}{Elle s'intéresse à l'histoire. --- L'histoire est intéressante.}

\begin{piege}[\en{*I am boring}]
Si vous dites \en{I'm boring}, vous dites « je suis quelqu'un d'ennuyeux ». Vous vouliez
dire « je m'ennuie » : \en{I'm bored}. Même piège avec \en{I'm excited} (j'ai hâte / je suis
enthousiaste) contre \en{I'm exciting} (je suis quelqu'un d'excitant), et
\en{I'm tired} contre \en{I'm tiring}.
\end{piege}

\subsection{Les 40 adjectifs de base}
\begin{tableau}[Adjectifs essentiels et leurs contraires]
\footnotesize
\begin{tabularx}{\linewidth}{@{}X X X X@{}}
\en{big} --- \en{small} (grand --- petit) & \en{long} --- \en{short} (long --- court) & \en{tall} --- \en{short} (grand --- petit, taille) & \en{old} --- \en{new / young} (vieux --- neuf / jeune) \\
\en{good} --- \en{bad} (bon --- mauvais) & \en{hot} --- \en{cold} (chaud --- froid) & \en{fast} --- \en{slow} (rapide --- lent) & \en{easy} --- \en{difficult / hard} (facile --- difficile) \\
\en{happy} --- \en{sad} (heureux --- triste) & \en{rich} --- \en{poor} (riche --- pauvre) & \en{clean} --- \en{dirty} (propre --- sale) & \en{full} --- \en{empty} (plein --- vide) \\
\en{strong} --- \en{weak} (fort --- faible) & \en{beautiful} --- \en{ugly} (beau --- laid) & \en{cheap} --- \en{expensive} (bon marché --- cher) & \en{open} --- \en{closed} (ouvert --- fermé) \\
\en{early} --- \en{late} (tôt --- tard) & \en{right} --- \en{wrong} (correct --- faux) & \en{true} --- \en{false} (vrai --- faux) & \en{safe} --- \en{dangerous} (sûr --- dangereux) \\
\en{heavy} --- \en{light} (lourd --- léger) & \en{wide} --- \en{narrow} (large --- étroit) & \en{high} --- \en{low} (haut --- bas) & \en{deep} --- \en{shallow} (profond --- peu profond) \\
\en{young} --- \en{old} (jeune --- âgé) & \en{kind} --- \en{mean} (gentil --- méchant) & \en{quiet} --- \en{noisy} (calme --- bruyant) & \en{free} --- \en{busy} (libre --- occupé) \\
\en{near} --- \en{far} (près --- loin) & \en{wet} --- \en{dry} (mouillé --- sec) & \en{dark} --- \en{bright} (sombre --- lumineux) & \en{interesting} --- \en{boring} (intéressant --- ennuyeux) \\
\end{tabularx}
\end{tableau}

\section{Unité 10 --- L'impératif, \en{can} et les formules de politesse}

\subsection{L'impératif}
\begin{regle}
Base verbale seule, sans sujet : \en{Come here!} \en{Don't touch that!} \en{Let's go!}
(Allons-y). Pour inclure quelqu'un : \en{Let's + base} ; négatif : \en{Let's not argue.}
\end{regle}
\ex{Sit down, please. Don't forget your bag. Let's start.}{Asseyez-vous, s'il vous plaît. N'oubliez pas votre sac. Commençons.}

\subsection{\en{can} : capacité, permission, possibilité}
\begin{tableau}[\en{can} --- invariable, jamais de \en{-s}, jamais de \en{to}]
\begin{tabularx}{\linewidth}{@{}p{3.6cm}p{4.4cm}X@{}}
\entete{Emploi} & \entete{Exemple} & \entete{Traduction} \tabularnewline
\hline
Capacité & \en{I can swim. She can't drive.} & Je sais nager. Elle ne sait pas conduire. \tabularnewline
Permission (informel) & \en{Can I open the window?} & Puis-je ouvrir la fenêtre ? \tabularnewline
Possibilité & \en{It can be dangerous.} & Ça peut être dangereux. \tabularnewline
Demande polie & \en{Can you help me, please?} & Pouvez-vous m'aider ? \tabularnewline
Plus poli : \en{could} & \en{Could you help me?} & Pourriez-vous m'aider ? \tabularnewline
Encore plus poli : \en{would you mind} & \en{Would you mind opening the door?} & Cela vous ennuierait-il d'ouvrir la porte ? \tabularnewline
\end{tabularx}
\end{tableau}

\begin{piege}[\en{*I can to speak English}]
Après un modal (\en{can, could, must, should, will, would, may, might}), le verbe est à la
base, \textbf{sans \en{to}} : \en{I can speak English.} \en{You must go.} \en{She will call you.}
\end{piege}

\subsection{Les formules de politesse (à apprendre comme des blocs)}
\begin{tableau}[La politesse anglaise, du plus direct au plus poli]
\footnotesize
\begin{tabularx}{\linewidth}{@{}X X@{}}
\en{Give me a coffee.} \textit{(ordre --- à éviter)} & \en{I want a coffee.} \textit{(trop direct)} \\
\en{Can I have a coffee, please?} \textit{(courant, parfait)} & \en{Could I have a coffee, please?} \textit{(plus poli)} \\
\en{I'd like a coffee, please.} \textit{(le plus naturel et sûr)} & \en{Would you mind passing me the sugar?} \textit{(très poli)} \\
\end{tabularx}
\end{tableau}

\begin{memo}[Les six blocs de politesse à automatiser]
\en{Could you\ldots, please?} \quad \en{Would you mind + \ldots ing?} \quad
\en{I'd like\ldots} \quad \en{May I\ldots?} \quad \en{Sorry to bother you, but\ldots} \quad
\en{Thanks a lot, I really appreciate it.}
\end{memo}

\section{Unité 11 --- Les prépositions de temps : \en{at}, \en{on}, \en{in}}

\begin{tableau}[La règle des trois prépositions de temps]
\begin{tabularx}{\linewidth}{@{}p{1.8cm}p{5.0cm}X@{}}
\entete{Mot} & \entete{S'emploie pour} & \entete{Exemples} \tabularnewline
\hline
\en{at} & Heures, moments précis, fêtes (sans \en{day}) & \en{at 7 o'clock, at noon, at midnight, at night, at the weekend (BrE), at Christmas, at the moment} \tabularnewline
\en{on} & Jours, dates, jours de fête & \en{on Monday, on 5 May, on my birthday, on Christmas Day, on Friday morning} \tabularnewline
\en{in} & Mois, années, saisons, siècles, durées, parties de la journée & \en{in May, in 2026, in summer, in the 21st century, in the morning, in two hours, in the past} \tabularnewline
\hline
\multicolumn{3}{@{}l}{\textbf{Sans préposition} : \en{today, tomorrow, yesterday, this week, last month, next year, every day, tonight}} \tabularnewline
\end{tabularx}
\end{tableau}

\begin{piege}[\en{*in Monday}, \en{*on the morning}, \en{*at July}]
Trois réflexes français à casser : \en{on Monday} (jamais \en{in}), \en{in the morning}
(mais \en{on Monday morning}), \en{at 8 o'clock} (jamais \en{in}). Et jamais de préposition
devant \en{next / last / this / every} : \en{I'll see you next week} (pas \en{*in next week}).
\end{piege}

\subsection{L'heure et les dates}
\begin{tableau}[Dire l'heure]
\begin{tabularx}{\linewidth}{@{}p{4.0cm}X@{}}
\entete{Heure} & \entete{Façons de la dire} \tabularnewline
\hline
8:00 & \en{It's eight o'clock.} \tabularnewline
8:05 & \en{It's eight oh five.} / \en{It's five past eight.} \tabularnewline
8:15 & \en{It's eight fifteen.} / \en{It's a quarter past eight.} \tabularnewline
8:30 & \en{It's eight thirty.} / \en{It's half past eight.} \tabularnewline
8:45 & \en{It's eight forty-five.} / \en{It's a quarter to nine.} \tabularnewline
8:50 & \en{It's eight fifty.} / \en{It's ten to nine.} \tabularnewline
\end{tabularx}
\end{tableau}

\begin{tableau}[Les dates]
\begin{tabularx}{\linewidth}{@{}p{4.0cm}X@{}}
\entete{Forme} & \entete{Écriture et prononciation} \tabularnewline
\hline
Britannique & \en{29/10/2026} = \en{29 October 2026} = \en{the twenty-ninth of October, twenty twenty-six} \tabularnewline
Américaine & \en{10/29/2026} = \en{October 29, 2026} = \en{October twenty-ninth} \tabularnewline
\textbf{Attention} & \en{10/11} = 10 novembre (BrE) ou 11 octobre (AmE) --- écrivez le mois en lettres dans les documents importants. \tabularnewline
\end{tabularx}
\end{tableau}

\subsection{Les nombres}
\begin{tableau}[Nombres de 0 à 1\,000\,000 --- pièges inclus]
\footnotesize
\begin{tabularx}{\linewidth}{@{}X X@{}}
0 \en{zero / oh} & 11 \en{eleven} \quad 12 \en{twelve} \\
1 \en{one} \quad 2 \en{two} \quad 3 \en{three} & 13 \en{thirteen} \quad 30 \en{thirty} \\
4 \en{four} \quad 5 \en{five} \quad 6 \en{six} & 14 \en{fourteen} \quad 40 \en{forty} (sans \en{u} !) \\
7 \en{seven} \quad 8 \en{eight} \quad 9 \en{nine} & 15 \en{fifteen} \quad 50 \en{fifty} \\
10 \en{ten} & 16 \en{sixteen} \quad 60 \en{sixty} \\
20 \en{twenty} \quad 21 \en{twenty-one} & 17 \en{seventeen} \quad 70 \en{seventy} \\
100 \en{a / one hundred} & 18 \en{eighteen} \quad 80 \en{eighty} \\
1\,000 \en{a / one thousand} & 19 \en{nineteen} \quad 90 \en{ninety} \\
1\,000\,000 \en{a million} & 2026 \en{twenty twenty-six} \tabularnewline
\end{tabularx}
\end{tableau}

\begin{astuce}[Les nombres ne prennent jamais de \en{-s}]
\en{two hundred people}, \en{three thousand cars}, \en{five million} --- jamais
\en{*two hundreds}. Mais : \en{hundreds of people}, \en{thousands of cars}
(avec \en{of}, le pluriel réapparaît).
\end{astuce}

\begin{exo}[11.1 --- Prépositions de temps]
Complétez avec \en{at}, \en{on}, \en{in} ou rien :
\begin{enumerate}
  \item \trou[1cm] Monday \item \trou[1cm] the morning \item \trou[1cm] 6 o'clock
  \item \trou[1cm] 2026 \item \trou[1cm] July \item \trou[1cm] tomorrow
  \item \trou[1cm] night \item \trou[1cm] my birthday \item \trou[1cm] last week
  \item \trou[1cm] the weekend \item \trou[1cm] two hours \item \trou[1cm] summer
\end{enumerate}
\end{exo}
\begin{corrige}[11.1]
1. \en{on} 2. \en{in} 3. \en{at} 4. \en{in} 5. \en{in} 6. rien 7. \en{at}
8. \en{on} 9. rien 10. \en{at} (BrE) / \en{on} (AmE) 11. \en{in} 12. \en{in}
\end{corrige}

\begin{exo}[11.2 --- Écrivez en toutes lettres]
\en{3:15, 7:45, 12:30, 14, 40, 19, 100, 2\,000, 15\,000, 05/12/2026 (BrE), 1987}.
\end{exo}
\begin{corrige}[11.2]
\en{three fifteen / a quarter past three ; seven forty-five / a quarter to eight ;
half past twelve ; fourteen ; forty ; nineteen ; a hundred ; two thousand ;
fifteen thousand ; the fifth of December twenty twenty-six ; nineteen eighty-seven}.
\end{corrige}

\begin{memo}[Révision express du chapitre 8]
\begin{enumerate}
  \item Sujet obligatoire toujours : \en{It is cold}, \en{It rains}, \en{There is a problem}.
  \item \en{be} pour l'âge, la faim, le froid, la peur, la raison.
  \item Présent simple : \en{-s} à la 3\textsuperscript{e} personne, \en{do/does} aux questions
  et négations, verbe de base après \en{do/does/did}.
  \item Article zéro pour les généralités, les repas, les langues, les sports, les pays.
  \item Verbes d'état jamais au \en{-ing}.
  \item Adjectif avant le nom, invariable.
  \item \en{at} + heure, \en{on} + jour, \en{in} + mois/année/période.
  \item Après un modal, base verbale sans \en{to}.
  \item Une seule négation par phrase.
  \item Question : auxiliaire + sujet + verbe.
\end{enumerate}
\end{memo}
